\documentclass[12pt,a4paper]{article}

% 패키지
\usepackage{graphicx}      % 이미지 삽입
\usepackage{amsmath}       % 수식
\usepackage{siunitx}       % 단위 표현
\usepackage{booktabs}      % 표 예쁘게
\usepackage{geometry}      
\geometry{margin=1in}

\title{Physics Lab Report \\ 
\large Measurement of Gravitational Acceleration}
\author{Your Name}
\date{\today}

\begin{document}

\maketitle
\tableofcontents
\newpage

\section{Introduction}

The purpose of this experiment is to measure the gravitational acceleration using a simple pendulum.
The theoretical formula for the period of a pendulum is:

\begin{equation}
T = 2\pi \sqrt{\frac{L}{g}}
\end{equation}

where:
\begin{itemize}
    \item $T$ is the period,
    \item $L$ is the length of the pendulum,
    \item $g$ is the gravitational acceleration.
\end{itemize}

\section{Method}

The length of the pendulum was measured using a ruler.
The time for 10 oscillations was recorded three times.

\section{Results}

\subsection{Measured Data}

\begin{table}[h]
\centering
\begin{tabular}{ccc}
\toprule
Trial & Time for 10 Oscillations (s) & Period (s) \\
\midrule
1 & 20.1 & 2.01 \\
2 & 19.8 & 1.98 \\
3 & 20.0 & 2.00 \\
\bottomrule
\end{tabular}
\caption{Measured oscillation times}
\end{table}

\subsection{Calculation}

From the formula:

\begin{equation}
g = \frac{4\pi^2 L}{T^2}
\end{equation}

Substituting the measured values gives:

\begin{equation}
g \approx 9.8 \, \si{m/s^2}
\end{equation}

\section{Discussion}

The measured value is close to the theoretical value of 
\SI{9.81}{m/s^2}. Small errors may have occurred due to reaction time.

\section{Conclusion}

The experiment successfully estimated the gravitational acceleration
with reasonable accuracy.

\end{document}
